% =============================================================================
% CITS4402 Computer Vision Project Report — Face Detection and Matching
% Preet Prasad — 23857493
%
% Compiles on Overleaf with pdfLaTeX (default).  Requires refs.bib in the
% same project for the bibliography step.  No external images, no custom
% fonts — drops into a blank Overleaf project as-is.
%
% Build order on a local TeX install:
%   pdflatex report.tex
%   bibtex   report
%   pdflatex report.tex
%   pdflatex report.tex
% Overleaf runs this automatically.
% =============================================================================
\documentclass[11pt,a4paper,titlepage]{article}

\usepackage[a4paper,top=1.8cm,bottom=1.8cm,left=2.0cm,right=2.0cm]{geometry}
\usepackage[T1]{fontenc}
\usepackage[utf8]{inputenc}
\usepackage{helvet}            % Arial-equivalent (Helvetica) sans-serif
\renewcommand{\familydefault}{\sfdefault}
\usepackage{graphicx}
\usepackage{titlesec}
\usepackage{enumitem}
\usepackage{tabularx}
\usepackage{booktabs}
\usepackage{array}
\usepackage{ragged2e}
\usepackage{xcolor}
\usepackage{amsmath}
\usepackage[numbers,sort&compress]{natbib}
\usepackage[hyphens]{url}
\usepackage{hyperref}
\hypersetup{colorlinks=true,urlcolor=blue,linkcolor=black,citecolor=black}

% Tighten section spacing so the body fits inside 3 pages (title page and
% references do not count toward the 3-page limit per the unit announcement).
\titlespacing*{\section}{0pt}{4pt}{1pt}
\titlespacing*{\subsection}{0pt}{3pt}{0pt}
\titleformat{\section}{\normalfont\bfseries\fontsize{12pt}{13pt}\selectfont}{\thesection.}{0.5em}{}
\titleformat{\subsection}{\normalfont\bfseries\fontsize{11pt}{12pt}\selectfont}{\thesubsection}{0.5em}{}

\setlist{nosep,leftmargin=1.3em,topsep=1pt,partopsep=0pt}
\setlength{\parindent}{0pt}
\setlength{\parskip}{1pt}

% Column types for the iteration tables.
\newcolumntype{V}{>{\RaggedRight\arraybackslash\small}p{0.12\linewidth}}
\newcolumntype{C}{>{\RaggedRight\arraybackslash\small}p{0.43\linewidth}}
\newcolumntype{O}{>{\RaggedRight\arraybackslash\small}p{0.38\linewidth}}

\renewcommand{\arraystretch}{1.0}

\begin{document}

% ── Title page (does not count toward the 3-page limit) ──────────────────────
\begin{titlepage}
  \centering
  \vspace*{2cm}
  \includegraphics[width=0.55\linewidth]{uwa_logo.png}\par
  \vspace{2.5cm}
  {\fontsize{18pt}{22pt}\bfseries\selectfont CITS4402 Computer Vision\par}
  \vspace{6pt}
  {\fontsize{16pt}{20pt}\bfseries\selectfont Project Report\par}
  \vspace{20pt}
  {\fontsize{14pt}{18pt}\bfseries\selectfont Face Detection and Matching\par}
  \vfill
  {\fontsize{13pt}{16pt}\selectfont Preet Prasad\par}
  \vspace{4pt}
  {\fontsize{13pt}{16pt}\selectfont Student Number: 23857493\par}
  \vspace*{3cm}
\end{titlepage}

% ── Body starts here — page 1 of the 3-page content budget ──────────────────
% ── 1. Introduction ──────────────────────────────────────────────────────────
\section{Introduction}
This report describes the face detection and identity clustering system implemented
for the CITS4402 project. The system processes 640$\times$480 colour images (1--4 faces),
detects faces, localises three landmarks (right eye, left eye, nose tip), aligns
each face to a 125$\times$125 crop, and clusters faces across a folder into unique
identities. The pipeline uses MediaPipe BlazeFace~\citep{bazarevsky2019blazeface}
and FaceLandmarker~\citep{lugaresi2019mediapipe} when available, with a
fully self-contained fallback built on OpenCV~\citep{bradski2000opencv} Haar
cascades~\citep{viola2001rapid} and the Hough circle transform~\citep{yuen1990hough}
covered in Lab~04. For clustering, 128-D dlib ResNet embeddings~\citep{king2009dlib}
(via the \texttt{face\_recognition} library) feed DBSCAN~\citep{ester1996density}.
The table below shows the headline measured outcomes from the naive starting
point to the final implementation; subsequent sections detail each task's
development.

\subsection{Baseline vs Final (measured on the 12 sample images)}
\noindent
\begin{tabularx}{\linewidth}{@{}p{0.22\linewidth} p{0.32\linewidth} X@{}}
\toprule
\textbf{Metric} & \textbf{Naive baseline} & \textbf{Final implementation} \\
\midrule
Face detection         & 11/12 images (4.jpg missed)   & 12/12 images, all counts correct \\
Landmark detection     & 61\% both-eyes (17/28 faces)  & 100\% both-eyes (26/26 faces) \\
Landmarks inside box   & not measured                  & 100\% RE, 100\% LE, 100\% nose \\
Alignment success      & 61\%, no similarity transform & 100\%, exact targets (0 px residual) \\
Identity clustering    & not implemented               & eps=0.38, 6 clusters from 26 faces \\
\bottomrule
\end{tabularx}

% ── 2. Face Detection ────────────────────────────────────────────────────────
\section{Face Detection (Task 2)}

\subsection{Development Iterations}
\noindent
\begin{tabularx}{\linewidth}{@{}V C O@{}}
\toprule
\textbf{Version} & \textbf{Change} & \textbf{Outcome} \\
\midrule
v1 (baseline) & Haar cascade~\citep{viola2001rapid} (\texttt{frontalface\_default}) only,
                no skin filter
              & 11/12 images; 4.jpg missed entirely \\
v2 & Added MediaPipe BlazeFace~\citep{bazarevsky2019blazeface} (short-range +
     full-range) as primary detector; Haar retained as fallback
   & 12/12 images, but false positives: 5 boxes on 11.jpg (small body region),
     1 face missed on 12.jpg (Ross failed cascade re-verification) \\
v3 & Added skin-coverage bypass for verification (skin $\geq$ 75\% AND box $\geq$ 60 px);
     cap result at 4 (spec maximum)
   & Fixed 12.jpg (Ross retained) and 11.jpg (small body box dropped); 12/12 correct \\
v4 & Replaced YCrCb-only skin filter with three-rule majority vote
     (YCrCb + HSV wrap-around + normalised-RGB)
   & More tolerant of unusual lighting/darker skin without losing FP rejection \\
v5 (final) & CLAHE~\citep{zuiderveld1994clahe} replaces \texttt{equalizeHist}
             for the Haar-fallback grayscale input
           & 12/12 maintained; better local detail for the Haar fallback path \\
\bottomrule
\end{tabularx}

\subsection{Final Method}
Primary: MediaPipe BlazeFace~\citep{bazarevsky2019blazeface} short-range (230~KB)
+ full-range (1.1~MB) TFLite models, run on the original and horizontally-flipped
image. Results merged with IoU-NMS (threshold 0.28, containment $>$ 72\%). Each
box expanded 20\% up, 10\% left/right, 15\% down to capture the full head. Capped
at 4 by ranking on confidence + skin-coverage score. Fallback when MediaPipe is
unavailable: \texttt{frontalface\_default} + \texttt{alt2} Haar
cascades~\citep{viola2001rapid}, 28\% upward expansion, ROI re-verification,
profile cascade as last resort.

\subsection{Skin-Colour Segmentation (Mandatory)}
Three-rule majority vote (Lab~05 technique): a pixel is skin when $\geq 2$ of
three rules agree --- YCrCb ($\text{Cr} \in [133, 183]$, $\text{Cb} \in [75, 135]$),
HSV ($\text{H} \in [0,28] \cup [160,179]$), normalised RGB
($\text{R} > \text{G}-5$, $\text{R} > \text{B}$, $\text{R} > 70$). Morphological
open + close (5$\times$5 ellipse) before the 20\% pixel-ratio threshold.

\subsection{Alternatives Considered}
\begin{itemize}
  \item dlib HOG face detector: lower recall on non-frontal faces vs BlazeFace.
  \item OpenCV YuNet: needs a separate ONNX file not bundled with \texttt{opencv-python}.
\end{itemize}

% ── 3. Landmark Detection ────────────────────────────────────────────────────
\section{Facial Landmark Detection (Task 3)}

\subsection{Development Iterations}
\noindent
\begin{tabularx}{\linewidth}{@{}V C O@{}}
\toprule
\textbf{Version} & \textbf{Change} & \textbf{Outcome} \\
\midrule
v1 (baseline) & \texttt{haarcascade\_eye.xml}~\citep{viola2001rapid} only
              & 61\% both-eyes (17/28); fails on tilted, glasses, occluded faces \\
v2 & Hough circle transform~\citep{yuen1990hough} with Lab~04 edge-overlap
     scoring on upper 58\% ROI; \texttt{eye\_tree\_eyeglasses} and directional
     eye cascades as chained fallbacks; nose tip via intensity minimum of
     Gaussian-blurred patch
   & 100\% both-eyes on all 26 faces \\
v3 & Added MediaPipe FaceLandmarker~\citep{lugaresi2019mediapipe} (iris
     landmarks 468/473, nose landmark 1) as primary path
   & 100\% maintained; sub-pixel iris precision \\
v4 & Per-eye position validation: reject MediaPipe landmarks where any single
     eye is outside the top 8--60\% of the face box, or eyes differ in y by
     $>$ 20\% face height
   & Caught a Chandler-in-7.jpg case where MediaPipe placed the right eye on
     the forehead, causing a badly tilted alignment \\
v5 (final) & Integral-projection profile + symmetry-mirror fallback tiers;
             CLAHE~\citep{zuiderveld1994clahe} preprocessing
           & 100\% even with MediaPipe disabled (verified by forcing the
             landmarker off) \\
\bottomrule
\end{tabularx}

\subsection{Final Method}
Seven-tier fallback chain (each tier only runs if the previous failed):
(1) MediaPipe FaceLandmarker on a padded ROI;
(2) Hough circles with Lab~04 edge-overlap scoring;
(3) \texttt{haarcascade\_eye\_tree\_eyeglasses};
(4) directional eye cascades;
(5) integral-projection profile (darkest row in the upper face);
(6) symmetry mirror if one eye is found;
(7) proportional geometric estimate.
Nose tip always from intensity minimum in the central 30--70\% $\times$ 50--72\%
region of the face ROI.

\subsection{Alternatives Considered}
\begin{itemize}
  \item dlib 68-point shape predictor: 68~MB model file, too large for submission.
\end{itemize}

% ── 4. Alignment ─────────────────────────────────────────────────────────────
\section{Facial Alignment (Task 4)}

\subsection{Development Iterations}
\noindent
\begin{tabularx}{\linewidth}{@{}V C O@{}}
\toprule
\textbf{Version} & \textbf{Change} & \textbf{Outcome} \\
\midrule
v1 (baseline) & \texttt{getRotationMatrix2D} for rotation, then separate crop
                + resize. Does not use a similarity transform; landmarks not
                at target positions
              & 61\% success (tied to landmark failures) \\
v2 & \texttt{cv2.estimateAffinePartial2D} with LMEDS on all three landmarks
   & 100\% success but visible tilt artefact when nose constraint is geometrically
     inconsistent with eyes (e.g., Monica in 7.jpg, eyes both at y=133 but face
     aligned with $\sim$20$^{\circ}$ rotation) \\
v3 & Analytical 2-point similarity transform from the two eye points only
   & 100\%, no tilt; eyes exact at (40,40) and (85,40); but nose drifted from
     (63,70) --- flagged by the lab facilitator \\
v4 & Least-squares similarity transform over all three landmarks (closed form
     via \texttt{np.linalg.lstsq} on the 6-equation 4-unknown system)
   & 100\%; reduced tilt vs LMEDS but still measured max 9.5~px residual
     because a 4-DOF model cannot exactly fit 3 arbitrary point pairs \\
v5 (final) & \texttt{cv2.getAffineTransform}~\citep{bradski2000opencv} with all
             three landmarks
           & 100\%; verified 0~px residual on all 26 faces --- every landmark
             lands precisely at its target coordinate \\
\bottomrule
\end{tabularx}

\subsection{Final Method}
\texttt{cv2.getAffineTransform} takes exactly 3 source$\rightarrow$target pairs
and returns the unique 6-DOF affine matrix that maps them exactly. A single
\texttt{cv2.warpAffine} call with output size (125, 125) simultaneously aligns,
crops, resizes, and (where needed) lightly stretches the face so that
$\text{right\_eye} \to (40, 40)$, $\text{left\_eye} \to (85, 40)$, and
$\text{nose\_tip} \to (63, 70)$ are all exact. \texttt{borderMode=BORDER\_REPLICATE}
fills out-of-bounds pixels, preventing black borders. Landmark dots on the
corner thumbnail are drawn at the fixed target coordinates and never move
between images.

% ── 5. Identity Clustering ───────────────────────────────────────────────────
\section{Identity Clustering (Task 5)}

\subsection{Development Iterations}
\noindent
\begin{tabularx}{\linewidth}{@{}V C O@{}}
\toprule
\textbf{Version} & \textbf{Change} & \textbf{Outcome} \\
\midrule
v1 (baseline) & Not implemented in the baseline code & --- \\
v2 & 128-D dlib ResNet embeddings~\citep{king2009dlib} via \texttt{face\_recognition}
     (explicitly permitted by spec); DBSCAN~\citep{ester1996density} via
     scikit-learn~\citep{pedregosa2011scikit} with \texttt{min\_samples=1} and
     Euclidean distance
   & Clusters identities meaningfully; see eps sweep below \\
v3 & Singleton-merge post-processing: merge isolated one-face clusters into
     the nearest compatible cluster if distance $<$ 0.515 AND no shared
     source image
   & Repairs cases where one image of a person has unusual pose/lighting \\
v4 & Upgraded the LBP fallback descriptor from basic 256-bin 5$\times$5 to
     LBP-u2~\citep{ojala2002multiresolution} 7$\times$7 multi-scale
     following \citet{ahonen2006face}: 59 uniform-pattern bins
     $\times$ 49 cells $\times$ 2 radii = 5782-D L2-normalised
   & LBP fallback peak silhouette~\citep{rousseeuw1987silhouettes}:
     0.04 $\to$ 0.11 ($\approx 3\times$ improvement) \\
\bottomrule
\end{tabularx}

\subsection{DBSCAN eps Sweep (final code, 26 faces)}
\noindent
\begin{tabular}{@{}lll@{}}
\toprule
\textbf{eps} & \textbf{Clusters} & \textbf{Silhouette}~\citep{rousseeuw1987silhouettes} \\
\midrule
0.35           & 10 & 0.1408 \\
\textbf{0.38}  & \textbf{6}  & \textbf{0.1021 $\leftarrow$ selected} \\
0.40           & 4  & 0.0954 \\
0.43           & 1  & 0.0000 \\
0.45           & 1  & 0.0000 \\
\bottomrule
\end{tabular}

\smallskip
eps=0.38 was selected because 6 clusters is closest to the 5-identity ground
truth (the 5 sample-image characters), with one slight over-split for a
character whose appearance varies significantly between images. eps=0.35 has
marginally higher silhouette (0.14) but over-splits to 10 clusters. eps=0.40
collapses two distinct identities. \texttt{min\_samples=1} ensures every face
joins a cluster (no noise points).

\subsection{Alternatives Considered}
\begin{itemize}
  \item k-means: cannot be used --- number of identities is unknown a priori.
  \item Agglomerative clustering: rejected because DBSCAN is explicitly named
        in the spec.
\end{itemize}

% ── 6. Final Results ─────────────────────────────────────────────────────────
\section{Final Results}
Measured on the 12 sample images (26 detected faces, all face counts match
ground truth):
\begin{itemize}
  \item Face detection: 12/12 images correct; avg detection $\approx$ 39~ms
        (MediaPipe path).
  \item Landmark detection: 100\% both eyes; 100\% RE, 100\% LE, 100\% nose
        inside the face box.
  \item Alignment: 100\% success; 0~px residual on every landmark; all outputs
        exactly 125$\times$125.
  \item Clustering: 6 DBSCAN clusters from 26 faces (eps=0.38, silhouette
        = 0.10).
\end{itemize}

\textbf{Known limitations.}
(1) MediaPipe model files ($\approx$ 4.9~MB total) must accompany the script
--- the Haar/Hough fallback covers their absence automatically.
(2) \texttt{face\_recognition} (dlib) must be pip-installed for best clustering
--- the LBP-u2 fallback works without it but with reduced discrimination.
(3) eps=0.38 over-splits one identity (6 clusters from 5 characters); eps=0.40
merges two distinct characters, so 6 is the better error.

% ── References (does not count toward the 3-page limit) ──────────────────────
\renewcommand{\bibsection}{\section*{References}}
\bibliographystyle{plainnat}
{\small\bibliography{refs}}

\end{document}
